\chapter{Graduate Attributes}{}
\makeatletter\@mkboth{}{Appendix}\makeatother
\label{appen:graduate_attributes}

\section{Complex Engineering Problem}
The project requires a first-principles understanding of CNNs, transfer learning and classification metrics (Chapter~\ref{chap:background}), and engagement with current research literature on remote sensing and image classification (Chapter~\ref{chap:literature}). The problem has no obvious solution. No study was found that classifies the operating state of treatment units from satellite images (Chapter~\ref{chap:literature}), so the problem is novel. The problem consists of several sub-problems: preparing and relabelling the dataset (Chapter~\ref{chap:systemmethods}), selecting and tuning architectures, handling class imbalance, and evaluating on known and new facilities (Chapter~\ref{chap:experimental_design}). It also involves non-technical issues, namely the cost of physical inspections and the compliance of South African WWTW (Chapter~\ref{chap:introduction}), and stakeholders outside engineering, such as plant operators and the annotators of the dataset.

\section{Problem Solving (GA 1)}
The problem was identified and formulated as a set of objectives in Chapter~\ref{chap:introduction}. The literature was reviewed in Chapter~\ref{chap:literature} to establish what had already been done and which methods were suitable. The problem was analysed by splitting it into two classification problems, one per component, and evaluating each on a chronological and a leave-one-facility-out split (Chapter~\ref{chap:experimental_design}). The conclusions in Chapter~\ref{chap:conclusion} are substantiated by the results in Chapter~\ref{chap:results}. Sustainable development is considered in that the envisaged system reduces the number of physical site visits, and supports earlier detection of failing treatment units, which reduces the discharge of insufficiently treated wastewater into the environment.

\section{Application of Scientific and Engineering Knowledge (GA 2)}
Mathematics and computing were applied in the formulation of the CNN building blocks and architectures (Section~\ref{sec:cnn_background}), the weighted cross-entropy loss (Section~\ref{sec:exp_loss}), the evaluation metrics (Sections~\ref{sec:metrics_background} and~\ref{sec:exp_metrics}) and Grad-CAM (Section~\ref{sec:gradcam_background}). Knowledge of the wastewater treatment process (Chapter~\ref{chap:background}) was required to understand which visual features indicate the state of an aerobic zone or clarifier, and to judge whether the Grad-CAM heatmaps identify these features (Section~\ref{sec:results_gradcam}).

\section{Engineering Design (GA 3)}
The classification system was designed in Chapter~\ref{chap:experimental_design}. Design choices include the preparation script (Section~\ref{sec:pipeline}), the data splits, the augmentations, the class weights, the hyperparameter search and the selection of a decision threshold for a target Not-OK recall of 95\% (Section~\ref{sec:exp_binary}). Each choice is motivated in the corresponding section. Public health and the environment are relevant, since a failing treatment unit releases insufficiently treated wastewater. The threshold is therefore chosen to favour recall of Not-OK components over specificity. Whole-life cost is relevant in that the system uses existing satellite images and requires no hardware on site, unlike on-site sensors, which are expensive to install and maintain (Chapter~\ref{chap:introduction}).

\section{Investigations, Experiments and Data Analysis (GA 4)}
Three experiments were designed and performed (Chapter~\ref{chap:experimental_design}), and their results analysed in Chapter~\ref{chap:results}. Critical thinking was applied to the data itself. Early inspection found labelling errors, so every clarifier image was reviewed, and not only the misclassified images, to avoid bias (Chapter~\ref{chap:systemmethods}). The data was split by date and by facility rather than at random, so that crops from the same satellite image do not appear in both the training and test sets. The Grad-CAM analysis was used to check whether correct classifications are based on relevant features, and found that some clarifier heatmaps lie on the rim of the tank (Section~\ref{sec:results_gradcam}). The limitations of the results are stated in Section~\ref{sec:conc_limitations}.

\section{Use of Engineering Tools (GA 5)}
The models were implemented and trained in PyTorch and torchvision, the hyperparameters were tuned with Optuna, the images were annotated with VIA2, and the results were analysed in Jupyter notebooks (Chapters~\ref{chap:systemmethods} and~\ref{chap:experimental_design}). Tools were also created: a preparation script that crops and masks the labelled regions (Section~\ref{sec:pipeline}), and a relabelling notebook built on Jupyter widgets. Limitations of the tools were recognised. Google Earth images contain only RGB bands and little spectral information (Chapter~\ref{chap:literature}), and the hyperparameter search was run for only one held-out facility per component due to its computational cost (Section~\ref{sec:exp_hyperparams}).

\section{Professional and Technical Communication (GA 6)}
This report presents the problem, method, results and conclusions in a structured way, with figures, diagrams and tables to support the text.

\section{Individual Work (GA 8)}
The project was carried out individually. This includes the literature study, the preparation and relabelling of the dataset, the implementation of the training and evaluation code, the experiments, the analysis of the results and the writing of the report. Weekly meetings with the supervisor were used to discuss progress and decisions, and the work was planned and managed according to the schedule in Appendix~\ref{appen:planning_schedule}.

\section{Independent Learning Ability (GA 9)}
CNN architectures, transfer learning, hyperparameter optimisation and Grad-CAM were learnt independently, mostly from the original papers and library documentation (Chapter~\ref{chap:background}). Feedback from the supervisor was discussed in the weekly meetings and incorporated into the experiments and the report. Uncertainty was present throughout, since it was not known beforehand whether the states could be determined from satellite images at all. Own limitations were recognised, for example in the labelling of ambiguous images, where the supervisor checked every proposed correction. The self-learning was sufficient to meet the objectives of the project.

\section{Sustainable Development}
The envisaged system replaces some physical site visits with remote monitoring, which reduces travel cost and emissions. It uses existing satellite images and requires no hardware on site, so no resources are spent on the installation and maintenance of sensors. Earlier detection of failing treatment units reduces the environmental impact of insufficiently treated effluent. Training CNNs consumes energy, which was limited by using pretrained models, early stopping and a single hyperparameter search per component.

\section{Critical Thinking}
Critical thinking was applied in questioning the labels, which led to the relabelling of the clarifiers, and in questioning the evaluation, which led to the leave-one-facility-out split to estimate performance at a new facility. The results were not accepted at face value: Grad-CAM was used to check what the models base their classifications on, and the limitations of the results are stated in Section~\ref{sec:conc_limitations}.

\section{Emerging Technologies}
Deep learning for image classification and remote sensing with satellite imagery are emerging technologies in the monitoring of wastewater treatment. Recent architectures such as ConvNeXt-Tiny were evaluated against older ones (Section~\ref{sec:exp_architectures}), and Grad-CAM was used to interpret the models. Higher-resolution or multispectral satellite data could further improve the classification (Section~\ref{sec:conc_future}).
